\begin{tikzpicture}
	\tikzmath{
		\widthpic = 8;
		\small = \widthpic / 400; %the filling rectangle in the nodes should be a bit submerged
		\l = 4; %distance between the nodes
	}

	%CENTRAL NODE PARAMETERS
	\tikzmath{
		\r = 0.8;
	}

	% draw central node
	\draw[fill=cw] (0, 0) circle (\r);

	\begin{scope}[rotate=45]
		%characterstics of 2nd node and properties
		\tikzmath{
			\d = 0.6; %thickness of the capillary
		}
		\tikzmath{
			\pb = sqrt(\r * \r - \d * \d / 4); % p border
			\pr = \pb - \small; %p length
		}
		\begin{scope}[shift={(\pr, 0)}]
			%default one
			\fill[cnw] (0, \d/2) rectangle (\l, -\d/2);



			%~ \begin{scope}
				%~ \tikzmath{
					%~ \pos1 = 0;
					%~ \pos2 = 0.2;
				%~ }
				%~ \fill[cnw] (\pos1 * \l, \d/2) rectangle (\pos2 * \l, -\d/2);
			%~ \end{scope}

			%draw a slob of rectangle
			\begin{scope}
				\tikzmath{
					\pos1 = 0;
					\pos2 = 0.2;
				}
				\fill[cdnw] (\pos1 * \l, \d/2) rectangle (\pos2 * \l, -\d/2);
			\end{scope}


			%drawing a meniscus
			\begin{scope}
				%SET MENISCUS POSITION
				\tikzmath{
					\meniscuspos = 0.03; %in 0-1 coordinates
					\meniscusrot = 180; %0=) 180=(
				}

				%drawing meniscus
				\begin{scope}[shift={(\meniscuspos*\l,0)}]
					\begin{scope}[rotate=\meniscusrot]
						\begin{scope}[shift={(-\d/4,0)}]
							\fill[cw] (0, \d / 2) rectangle (\d/2 + \small, -\d / 2);
							\fill[cdnw] (0, -\d/2) arc(-90:90:\d /2) -- cycle;
							\draw[] (0, -\d/2) arc(-90:90:\d /2);
						\end{scope}
					\end{scope}
				\end{scope}
			\end{scope}
			%draw the main outer borders
			\draw[] (0, \d/2) -- (\l, \d/2);
			\draw[] (0, -\d/2) -- (\l, -\d/2);

			%\draw[->, thick] (0.1*\l, 0) -- (0.5*\l, 0);
			%\node[] at (0.25*\l, -0.6) {$Q_1$};
		\end{scope}
	\end{scope}


	\begin{scope}[rotate=135]
		%characterstics of 2nd node and properties
		\tikzmath{
			\d = 0.4; %thickness of the capillary
		}
		\tikzmath{
			\pb = sqrt(\r * \r - \d * \d / 4); % p border
			\pr = \pb - \small; %p length
		}
		\begin{scope}[shift={(\pr, 0)}]
			%default one
			\fill[cnw] (0, \d/2) rectangle (\l, -\d/2);

			%draw a slob of rectangle
			\begin{scope}
				\tikzmath{
					\pos1 = 0;
					\pos2 = 0.6;
				}
				\fill[cdnw] (\pos1 * \l, \d/2) rectangle (\pos2 * \l, -\d/2);
			\end{scope}

			%drawing a meniscus
			\begin{scope}
				%SET MENISCUS POSITION
				\tikzmath{
					\meniscuspos = 0.03; %in 0-1 coordinates
					\meniscusrot = 180; %0=) 180=(
				}
				%drawing meniscus
				\begin{scope}[shift={(\meniscuspos*\l,0)}]
					\begin{scope}[rotate=\meniscusrot]
						\begin{scope}[shift={(-\d/4,0)}]
							\fill[cw] (0, \d / 2) rectangle (\d/2 + \small, -\d / 2);
							\fill[cdnw] (0, -\d/2) arc(-90:90:\d /2) -- cycle;
							\draw[] (0, -\d/2) arc(-90:90:\d /2);
						\end{scope}
					\end{scope}
				\end{scope}
			\end{scope}

			%draw the main outer borders
			\draw[] (0, \d/2) -- (\l, \d/2);
			\draw[] (0, -\d/2) -- (\l, -\d/2);
			%\draw[->, thick] (0.1*\l, 0) -- (0.2*\l, 0);
			%\node[] at (0.25*\l, -0.6) {$Q_2$};
		\end{scope}
	\end{scope}

	%tube
	\begin{scope}[rotate=-45]
		%characterstics of 2nd node and properties
		\tikzmath{
			\d = 0.8; %thickness of the capillary
		}
		\tikzmath{
			\pb = sqrt(\r * \r - \d * \d / 4); % p border
			\pr = \pb - \small; %p length
		}
		\begin{scope}[shift={(\pr, 0)}]
			%default one
			\fill[cnw] (0, \d/2) rectangle (\l, -\d/2);

			%drawing a meniscus
			\begin{scope}
				%SET MENISCUS POSITION
				\tikzmath{
					\meniscuspos = 0.03; %in 0-1 coordinates
					\meniscusrot = 180; %0=) 180=(
				}
				%drawing meniscus
				\begin{scope}[shift={(\meniscuspos*\l,0)}]
					\begin{scope}[rotate=\meniscusrot]
						\begin{scope}[shift={(-\d/4,0)}]
							\fill[cw] (0, \d / 2) rectangle (\d/2 + \small, -\d / 2);
							\fill[cnw] (0, -\d/2) arc(-90:90:\d /2) -- cycle;
							\draw[] (0, -\d/2) arc(-90:90:\d /2);
						\end{scope}
					\end{scope}
				\end{scope}
			\end{scope}


			%draw the main outer borders
			\draw[] (0, \d/2) -- (\l, \d/2);
			\draw[] (0, -\d/2) -- (\l, -\d/2);
			%\draw[<-, thick] (0.1*\l, 0) -- (0.4*\l, 0);
			%\node[] at (0.25*\l, -0.6) {$Q_3$};
		\end{scope}
	\end{scope}




	\begin{scope}[rotate=-135]
		%characterstics of 2nd node and properties
		\tikzmath{
			\d = 0.8; %thickness of the capillary
		}
		\tikzmath{
			\pb = sqrt(\r * \r - \d * \d / 4); % p border
			\pr = \pb - \small; %p length
		}
		\begin{scope}[shift={(\pr, 0)}]
			%default one
			\fill[cw] (0, \d/2) rectangle (\l, -\d/2);

			%draw the main outer borders
			\draw[] (0, \d/2) -- (\l, \d/2);
			\draw[] (0, -\d/2) -- (\l, -\d/2);
			%\draw[<-, thick] (0.1*\l, 0) -- (0.4*\l, 0);
			%\node[] at (0.25*\l, -0.8) {$Q_4$};
		\end{scope}
	\end{scope}

\end{tikzpicture}
